\par\Needspace{8\baselineskip}\bigskip
\pdfbookmark[3]{Exposition, 8}{pallas-II-article8}
\begin{center}8.\end{center}
\noindent\hspace*{\parindent}In this investigation we shall use the same symbols that we introduced in the first section, both for the elements proper and for the other auxiliary quantities. Thus, for example, \(A\) will denote the now variable value of \(\dfrac{y\mathrm{d}z-z\mathrm{d}y}{\mathrm{d}t}\). We see that all the finite equations relating these quantities \(A,B,C,F,G,H,v,E\) etc., the elements, and the coordinates \(x,y,z,r\) together with their differentials \(\dfrac{\mathrm{d}x}{\mathrm{d}t}\mathpunct{\raisebox{.45ex}{,}}\dfrac{\mathrm{d}y}{\mathrm{d}t}\mathpunct{\raisebox{.45ex}{,}}\dfrac{\mathrm{d}z}{\mathrm{d}t}\mathpunct{\raisebox{.45ex}{,}}\dfrac{\mathrm{d}r}{\mathrm{d}t}\) will be independent of this variability. We may therefore differentiate any finite equation relating the coordinates \(x,y,z,r\), the elements, and the other quantities that are constant under the elliptic hypothesis and may also be included under the name of elements, considering only the differentials of the coordinates. Subtracting this result from the complete differential leaves the differential of the given equation when the elements alone are treated as variable. This latter differentiation is therefore evidently legitimate as well. This remarkable theorem may sometimes shorten the calculations, although it is clear that we can always dispense with it.
